% Shared student/instructor body; compile a wrapper, not this file directly.
\providecommand{\dd}{\,\mathrm{d}}
\AtBeginDocument{%
\setlength{\abovedisplayskip}{4pt plus 1pt minus 1pt}
\setlength{\belowdisplayskip}{4pt plus 1pt minus 1pt}
\setlength{\abovedisplayshortskip}{3pt plus 1pt}
\setlength{\belowdisplayshortskip}{3pt plus 1pt}
\setlength{\jot}{2pt}
}
\tcbset{handout base/.append style={before skip=2.5pt,after skip=2.5pt,before upper={\RaggedRight}}}

\HandoutSetup{NE 630}{14}{The $k$-Eigenvalue Problem}
  {Review Lessons 12--13; bring the multigroup balance equations}
\begin{document}
\MakeHandoutTitle
\textcolor{KSUDeepPurple}{\textbf{Primary objective.}}
Compute the multiplication factor and multigroup spectrum of an
infinite, homogeneous system, and distinguish eigenvalue from fixed-source solutions.
\par\vspace{5pt}

\begin{handoutcolumns}
\HandoutSection{1}{Keep losses and production separate}
Recall $\Sigma_{r,g}=\Sigma_{t,g}-\Sigma_{s,g\gets g}$.
For compact algebra, define
\[
 f_g=(\nu\Sigma_f)_g,
 \qquad \mathcal P=\sum_{g'=1}^G f_{g'}\phi_{g'}.
\]
Then the physical steady-state balance is
\[
 \Sigma_{r,g}\phi_g=
 \sum_{g'\ne g}\Sigma_{s,g\gets g'}\phi_{g'}
       +\chi_g\mathcal P+s_g.
\]
Here $f_g$ is a production cross section in $\mathrm{cm^{-1}}$,
while $\mathcal P$ is a neutron-production rate density.
For the supplied data, $f_g=\bar\nu_g\Sigma_{f,g}$ and
$\Sigma_{a,g}=\Sigma_{\gamma,g}+\Sigma_{f,g}$; $\bar\nu_g$ is fission-rate weighted.

\HandoutSection{2}{A multiplying two-group system}
Use bounds $(E_0,E_1,E_2)=(10^7,1,10^{-3})$ eV.
Group 1 is nonthermal and group 2 is thermal.
All macroscopic cross sections below are in $\mathrm{cm^{-1}}$.
\begin{center}
\renewcommand{\arraystretch}{1.32}
\begin{tabular}{@{}lcc@{}}
\toprule
 & $g=1$ & $g=2$\\
\midrule
$\Sigma_{\gamma,g}$ & $2.04\times10^{-3}$ & $1.91\times10^{-2}$\\
$\Sigma_{f,g}$ & $2.96\times10^{-4}$ & $1.38\times10^{-2}$\\
$\bar\nu_g$ & $2.49$ & $2.43$\\
$\Sigma_{s,1\gets g}$ & $7.21\times10^{-1}$ & $0$\\
$\Sigma_{s,2\gets g}$ & $4.27\times10^{-2}$ & $2.72$\\
$v_g$ ($\mathrm{m/s}$) & $10^5$ & $2\times10^3$\\
\bottomrule
\end{tabular}
\end{center}
Assume no leakage or out-of-range losses, no upscatter, all
fission neutrons born in group 1, and a group-1 external source:
\[
 (\chi_1,\chi_2)=(1,0),\qquad(s_1,s_2)=(s_1,0).
\]
For numerical fluxes, use $\boxed{s_1=1\unit{cm^{-3}\,s^{-1}}}$.
\begin{boardbox}{Reduce the data before solving}
Set $d=\Sigma_{s,2\gets1}=0.0427\unit{cm^{-1}}$. Then
\begin{align*}
 \Sigma_{a,1}&=0.002336,&\Sigma_{a,2}&=0.0329,\\
 \Sigma_{r,1}&=\Sigma_{a,1}+d=0.045036,\\
 f_1&=0.00073704,&f_2&=0.033534.
\end{align*}
All of these coefficients have units $\mathrm{cm^{-1}}$.
Why is $\Sigma_{r,2}=\Sigma_{a,2}$ here?
\RuledLines{2}
\end{boardbox}

\switchcolumn
\RaggedRight
\HandoutSection{3}{Solve the fixed-source problem}
The two balances are
\begin{align*}
 \Sigma_{r,1}\phi_1&=f_1\phi_1+f_2\phi_2+s_1,\\
 \Sigma_{a,2}\phi_2&=d\phi_1.
\end{align*}
\begin{boardbox}{First the shape, then the magnitude}
The thermal equation determines the flux ratio:
\[
 R=\frac{\phi_2}{\phi_1}=\frac{d}{\Sigma_{a,2}}
 =\MathBlank[1.02in]{1.29787234}.
\]
Insert $\phi_2=R\phi_1$ into the first balance:
\[
 \boxed{\phi_1=\frac{s_1}{\Sigma_{r,1}-f_1-f_2R}},
 \qquad \boxed{\phi_2=R\phi_1}.
\]
The denominator is $0.00077610894\unit{cm^{-1}}$.
For the stated unit source,
\begin{align*}
 \phi_1&=\MathBlank[1.15in]{1288.5}\unit{cm^{-2}\,s^{-1}},\\
 \phi_2&=\MathBlank[1.15in]{1672.3}\unit{cm^{-2}\,s^{-1}}.
\end{align*}
\end{boardbox}
For another source strength, the fluxes scale in direct proportion
to $s_1$ while the material data remain fixed.

\HandoutSection{4}{Audit the solution}
Adding the group balances cancels intergroup scattering:
\[
 \boxed{\Sigma_{a,1}\phi_1+\Sigma_{a,2}\phi_2
       =f_1\phi_1+f_2\phi_2+s_1}.
\]
Using unrounded fluxes, the rates in $\mathrm{cm^{-3}\,s^{-1}}$ are
\begin{align*}
 R_a&=\MathBlank[1.20in]{58.02794},\\
 \mathcal P&=\MathBlank[1.20in]{57.02794},\\
 R_a-\mathcal P&=\MathBlank[1.20in]{1=s_1}.
\end{align*}
\begin{checkpoint}[Why were speeds supplied?]
Speeds do not enter these steady-state flux balances. Using
$v_g$ as a representative group speed gives
$n_g\approx\phi_g/v_g$. Convert the tabulated speeds to
$\mathrm{cm/s}$, then compare $n_2/n_1$ with $\phi_2/\phi_1$.
\[
 \frac{n_2}{n_1}\approx
 \frac{\phi_2}{\phi_1}\frac{v_1}{v_2}
 =\MathBlank[1.00in]{64.8936}.
\]
\RuledLines{2}
\end{checkpoint}
\end{handoutcolumns}

\newpage
\begin{handoutcolumns}
\HandoutSection{5}{Remove the source; introduce $k$}
With $s_1=0$, a nonzero \emph{physical steady-state} solution
requires exactly critical material. Instead, define an eigenproblem
by dividing \textbf{only fission production} by $k_{\infty}$:
\begin{align*}
 \Sigma_{r,1}\phi_1&=
      \frac{1}{k_{\infty}}(f_1\phi_1+f_2\phi_2),\\
 \Sigma_{a,2}\phi_2&=d\phi_1.
\end{align*}
The second equation is unchanged, so $R=d/\Sigma_{a,2}$.
\begin{boardbox}{Derive the eigenvalue}
Divide the first equation by $\phi_1>0$:
\[
 k_{\infty}\Sigma_{r,1}=f_1+f_2R.
\]
\[
 \boxed{k_{\infty}=\frac{f_1+f_2d/\Sigma_{a,2}}{\Sigma_{r,1}}}
 =\MathBlank[1.00in]{0.98276692}.
\]
The unmodified, source-free system is therefore
\RevealBlank[1.15in]{subcritical}.
\end{boardbox}
The $1/k_{\infty}$ factor is mathematical normalization of
production, not a physical source or a time-growth rate.

\HandoutSection{6}{Write the matrix eigenproblem}
Define the loss-minus-scattering and fission matrices:
\[
 \mathbf M=
 \begin{bmatrix}\Sigma_{r,1}&0\\-d&\Sigma_{a,2}\end{bmatrix},
 \qquad
 \mathbf F=
 \begin{bmatrix}f_1&f_2\\0&0\end{bmatrix}.
\]
\begin{fixedbox}{Fixed source versus eigenvalue}
\begin{align*}
 (\mathbf M-\mathbf F)\symbfit{\phi}&=\mathbf s
           &&\text{(fixed source)},\\
 \mathbf M\symbfit{\phi}&=\frac1{k_{\infty}}
                   \mathbf F\symbfit{\phi}
           &&\text{(eigenvalue)}.
\end{align*}
When $\mathbf M$ is invertible,
\[
 \boxed{\mathbf M^{-1}\mathbf F\symbfit{\phi}
       =k_{\infty}\symbfit{\phi}}.
\]
\end{fixedbox}
For $G$ groups the same structure uses
\[
 M_{gg'}=\Sigma_{t,g}\delta_{gg'}-\Sigma_{s,g\gets g'},
 \qquad F_{gg'}=\chi_g f_{g'}.
\]
\begin{checkpoint}
Why is the off-diagonal entry $-d$ in $\mathbf M$ negative,
although the scattering cross section itself is positive?
\RuledLines{2}
\end{checkpoint}

\switchcolumn
\RaggedRight
\HandoutSection{7}{Normalize a shape, not a power}
An eigenvector has arbitrary magnitude. One convenient convention is
\[
 \widehat\phi_g=\frac{\phi_g}{\phi_1+\phi_2},
 \qquad \widehat\phi_1+\widehat\phi_2=1.
\]
\begin{boardbox}{Report the eigen-spectrum}
\[
 \widehat\phi_1=\frac{1}{1+R}
   =\MathBlank[0.95in]{0.435185},
\]
\[
 \widehat\phi_2=\frac{R}{1+R}
   =\MathBlank[0.95in]{0.564815}.
\]
These are dimensionless fractions of total flux, not neutron
number fractions. What additional information would set a
physical flux magnitude?
\RuledLines{2}
\end{boardbox}
For complete energy coverage, summing the eigen-equations gives
\[
 \boxed{k_{\infty}=
 \frac{\sum_g f_g\phi_g}{\sum_g\Sigma_{a,g}\phi_g}}.
\]
Evaluate this with the \emph{eigen-spectrum}; it recovers the
one-group production-to-absorption ratio from Lesson 12.

\HandoutSection{8}{Connect to source multiplication}
For this model, the external source and fission births both enter
group 1. Therefore the fixed-source and eigenvalue problems have
the same flux ratio $R$.
Using $f_1+f_2R=k_{\infty}\Sigma_{r,1}$ in Section 3 yields
\[
 \boxed{\phi_1=\frac{s_1}{\Sigma_{r,1}(1-k_{\infty})}},
 \qquad \phi_2=R\phi_1.
\]
If fission \emph{production} is suppressed while keeping the
same loss and scattering coefficients, the group-1 flux is
$\phi_1^{(0)}=s_1/\Sigma_{r,1}$. Hence
\[
 \frac{\phi_1}{\phi_1^{(0)}}=\frac1{1-k_{\infty}}
 =\MathBlank[1.05in]{58.02794}.
\]
\begin{checkpoint}[Approach criticality from below]
At fixed $s_1>0$, what happens as $k_{\infty}\to1^{-}$?
Why is the negative algebraic flux obtained for $k_{\infty}>1$
not a physical steady-state solution?
\RuledLines{2}
\end{checkpoint}
\begin{takeawaybox}
A fixed source sets the flux magnitude in a subcritical system.
The eigenproblem sets $k_{\infty}$ and a relative spectrum, not
power. The simple $1/(1-k_{\infty})$ factor here relies on this
source shape and fixed two-group model; it is not universal.
\end{takeawaybox}
\end{handoutcolumns}
\end{document}
